\section{Reproducibility and the living benchmark}
\label{sec:repro}

The harness (\texttt{research/harness}, package \texttt{jigbench}) takes a YAML configuration, plans the
trials, starts any model servers the conditions need, runs the trials and appends one JSON record per
finished trial. Runs resume: finished trial identifiers are skipped, and an interrupted trial leaves no
record. Each run session records the Jig commit (and a hash of any uncommitted core changes), the
hardware, and for each endpoint the served model file, its quantisation, the llama.cpp build, the
context size and the number of slots, together with the full configuration and seeds.
\texttt{jigbench report} regenerates every table and figure in Section~\ref{sec:results} from the
records, and \texttt{jigbench submit-results} validates a run against the results schema and packages
it, with a checksum, for contribution by pull request.

Heavy runs follow a compute policy suited to a shared personal machine: they start only between 01:00
and 07:00 local time, or after 30 minutes without user input, and only when the GPU is not busy with
other work; they stop within seconds of the owner returning, and resume the next night.
